% Section B.1: the answers to the parts of lectures/L01/appendix/c_exercises.md that are not code.
% Every testbench message quoted was produced by a design built to Chapters 1 to 5 and broken in
% exactly the way the part describes.

\section{Chapter 1: The Register Package and the Peripheral Top}
\label{sol:c1}
Exercise~\ref{c1:ex:2} is the build of \code{uart_top}, checked by a clean analysis now and by
\code{uart_top_tb} from \chapref{5} on; only its closing question, part e), is answered here.

\solution{c1:ex:1}{\code{uart_def}}
\ans{a} The register indices are \code{REG_STATUS} 0, \code{REG_CTRL} 1, \code{REG_BAUD_DIV} 2,
\code{REG_TX_DATA} 3, \code{REG_RX_DATA} 4, \code{REG_RX_POP} 5 and \code{REG_ERR_FLAGS} 6, each the
byte offset divided by four. The \code{STATUS} bits are \code{ST_TX_READY} 0, \code{ST_RX_VALID} 1,
\code{ST_ERROR} 2 and \code{ST_TX_IDLE} 3; the \code{CTRL} bits \code{CT_ENABLE} 0,
\code{CT_PARITY_LO} 1, \code{CT_PARITY_HI} 2, \code{CT_STOP} 3, \code{CT_RX_IRQ} 4 and
\code{CT_TX_IRQ} 5; the \code{ERROR_FLAGS} bits \code{ER_FRAMING} 0, \code{ER_PARITY} 1 and
\code{ER_OVERRUN} 2.

\ans{b} In VHDL the position indexes the vector directly:

\begin{vhdlcode}
if (status(ST_RX_VALID) = '1') then
\end{vhdlcode}

\noindent In C++ the same number is shifted into a mask at the use site:
\code{(status & (1U << status::RX_VALID)) != 0U}. The protocol speaks of ``\code{STATUS} bit 1'',
and a position is that number and nothing else, so both files can be checked against the protocol
by reading it off, digit for digit. A mask would be a derived value, \code{2} here and \code{4} for
bit 2, which the reader has to recompute to compare, and which only means something for one word
width. Storing positions also keeps the two languages identical, because the only difference between
them is the operator that uses the constant, and that operator is written where it is used.

\ans{c} \textbf{In no testbench at all.} Every testbench that touches those two registers names them
through \code{uart_def}, and so does the register bank. \code{uart_regs_tb} writes \code{RX_POP} to
``\code{REG_RX_POP}'' and the bank decodes ``\code{REG_RX_POP}'', so after the swap both sides mean
6 and the bench still passes; the same holds for \code{ERROR_FLAGS} at 5. \code{uart_top_tb} never
reads or writes index 5 or 6. Both were run with the swapped package, and both report that they
passed. A constant that is wrong in the one place everything reads it from is wrong consistently,
and a consistent error is invisible to any check built from the same source.

The C++ side is a separate transcription, and it still says \code{RX_POP} is 5 and
\code{ERROR_FLAGS} is 6, so the host suites in Chapters~8 and~9 pass too. The swap surfaces only
when the two transcriptions finally meet, on the bench in \chapref{10}. There the driver's
\code{read()} ends with \code{writeReg(RX_POP, 1)}, which the peripheral now takes as an
\code{ERROR_FLAGS} write: it sets the framing flag, so \code{STATUS} bit 2 comes up, and it never
advances the RX FIFO, so \code{RX_VALID} stays set and every later \code{read()} returns the same
byte. \code{clearErrors()} writes index 6, which the peripheral now takes as a pop, and silently
throws a received byte away.

Rung c of the ladder can still pass, because it sends one byte and reads one byte back, and that
first read is correct; it shows the fault only if the bring-up program looks at \code{STATUS} after
the read, or sends a second byte. Rung d shows it for certain: the first character typed comes back
correctly and so does every read after it, the same character each time. At rung e,
\code{EchoNode} echoes the first character typed without end.

\solution{c1:ex:2}{\code{uart_top}}
\ans{e} With only the provided files and your skeleton in \file{hw/}, \code{make build-vhdl} analyzes
five modules, runs the two transport testbenches, and skips the rest:

\begin{console}
==> uart_top_tb (skipped: no baud_gen.vhd sync.vhd fifo.vhd uart_tx.vhd uart_rx.vhd uart_regs.vhd yet)
2 testbench(es) run, 7 skipped.
\end{console}

\noindent \code{baud_gen} and \code{uart_tx} arrive in \chapref{2}, \code{sync} and \code{uart_rx} in
\chapref{3}, \code{fifo} in \chapref{4} and \code{uart_regs} in \chapref{5}. The one that
\code{uart_top} never instantiates is \code{fifo}: \code{uart_regs} owns both FIFOs. The build needs
it anyway because elaboration needs every entity in the hierarchy, not only the top's direct
children, and \code{uart_regs} cannot be analyzed until the entity it instantiates twice is already
in the \code{work} library.

\solution{c1:ex:3}{What a clean analysis does not prove}
\ans{a} \code{ghdl -a} accepts it. Both ports are \code{in std_logic}, so the entity is as legal
with them swapped as without, and the architecture refers to them by name, so it still analyzes.
Nothing binds to the entity positionally until \code{uart_top_tb} runs in \chapref{5}, and there
the first check fails: the bench's clock now arrives on the peripheral's \code{mosi} and its data on
\code{sclk}, so the first write never lands and the read-back reports

\begin{console}
uart_top_tb: BAUD_DIV read-back mismatch, got 0x0000!
\end{console}

\noindent four chapters after the mistake was made, in a message about \code{BAUD_DIV}. That is the
cost of a positional contract: the order of the ports is checked by nothing except a simulation of
the whole system, and when that simulation fails it points at the first symptom rather than the
cause. A named association would have caught it at analysis; the course binds positionally so that
the order in each interface table is something you have to get right.

\ans{b} All zeros. Every read transaction clocks out \code{0x00000000}, whatever register it
addresses, because the bridge latches \code{reg_rdata} one cycle after the command byte and shifts
it out on \code{MISO}. It is harmless until \chapref{5} because nothing reads it: the only testbench
that drives \code{uart_top} is skipped until then, and no register exists yet to be misreported.
Without the placeholder the vector is undriven, so every bit is \code{'U'}, and the master samples
\code{'U'} on all 32 data bits of every read. Both were simulated: the skeleton with the placeholder
returned \code{00000000000000000000000000000000} for \code{BAUD_DIV} and for \code{STATUS}, and the
skeleton without it returned 32 \code{'U'}s for each. The placeholder makes the answer a defined
value; it does not make it a right one.

\ans{c} A release that lands inside the recovery and removal window of the flip-flops around a
clock edge leaves some of them out of reset on that edge and others on the next, and any of them
may go metastable on the way. The design then starts half-reset: a state machine whose state
register has left reset while its outputs have not, a FIFO whose pointers start one cycle apart,
the SPI slave's edge detector seeing a false edge on a synchronizer stage that left reset a cycle
early. None of those states is one the logic was designed to leave, and nothing brings it back
short of another reset.

\code{sync} cannot do this job because of its port list. It takes \code{reset_s2_n} as an input:
it needs a clean reset in order to reset itself, so it cannot be the thing that produces one. A reset
synchronizer needs the raw \code{reset_n} on its flops' asynchronous clear, so that the assertion
reaches them with no clock at all, and a constant \code{'1'} clocked through the chain, so that the
release arrives a clean two edges later. \code{sync} has no port for either: its only asynchronous
input is the already-synchronized reset, and its data input is what it synchronizes.

\ans{d} Because the top's instantiation line is the block's contract. \code{uart_top} declares the
signals it connects and binds them by position, so writing that line fixes how many ports the
block has, in which order, of which types and in which directions, before the block exists. For
\code{baud_gen} to bind without the top changing, its entity must declare exactly
\code{clock} and \code{reset_s2_n} (\code{in std_logic}), then \code{div}
(\code{in natural range 1 to 65535}), then \code{tick} (\code{out std_logic}), in that order, which is
also the order \code{baud_gen_tb} binds. The names inside may be anything; the positions and the
types may not.
